\section{Introduction}
The service mesh has become a foundational mechanism for traffic management, observability, and policy enforcement in microservice deployments, but its original sidecar-centric realization has imposed measurable compute and latency overhead on production systems. This overhead, often described as the sidecar tax, emerges from duplicated proxies, additional context switches, memory footprints per pod, and increased request-path complexity under high fan-out communication graphs \cite{envoyDataPlaneOverhead,serviceMeshPerfStudies}. In multi-tenant clusters where economic efficiency and strict performance isolation must coexist, the sidecar tax is not merely a cost issue; it is an architectural tax that constrains packing density and distorts capacity planning.

Sidecarless approaches seek to reduce this tax by consolidating data-plane functionality into shared node-level agents and kernel-resident packet paths, frequently leveraging eBPF and XDP for low-overhead traffic processing. While this shift can reduce average latency, tail behavior remains the dominant concern for user-facing reliability. In this proposal, p99 tail-latency is adopted as the primary performance metric because it captures the rare but consequential slow requests that directly affect service-level objectives and user-perceived quality. The central argument is that mean latency improvements are insufficient if p99 latency exhibits inter-tenant leakage under contention; therefore, any claim of sidecarless superiority must be validated against isolation quality at the tail.
